\chapter{Validation, qualité et TDD}

\section{Great Expectations}
Great Expectations formalise des règles de qualité sur les données : colonnes attendues, types, valeurs admissibles et cohérence globale \cite{great_expectations_validation}. Le projet génère deux artefacts :
\begin{itemize}
  \item \path{reports/great_expectations_suite.json} ;
  \item \path{reports/great_expectations_validation.json}.
\end{itemize}

\begin{lstlisting}
.\.venv\Scripts\python.exe -m cyberguard_ml.pipeline.great_expectations_check
Get-Content .\reports\great_expectations_validation.json
\end{lstlisting}

\begin{table}[H]
\centering
\caption{Résultat Great Expectations}
\label{tab:ge_result}
\begin{tabularx}{0.75\linewidth}{@{}lC@{}}
\toprule
\textbf{Indicateur} & \textbf{Valeur}\\
\midrule
Validation & Réussie\\
Nombre de lignes & 3 000\\
Taux d'attaque & 84,33 \%\\
Erreurs & 0\\
Avertissements & 0\\
\bottomrule
\end{tabularx}
\end{table}

\section{Qualité automatique}
La chaîne qualité vérifie le code avant publication :
\begin{itemize}
  \item \tooltag{Black} : formatage ;
  \item \tooltag{isort} : imports ;
  \item \tooltag{Ruff} : linting moderne ;
  \item \tooltag{MyPy} : typage statique ;
  \item \tooltag{Pytest} : tests unitaires et d'intégration.
\end{itemize}

\screenshotfigure[0.56\textheight]{05_quality_tests.png}{Tests et contrôles qualité.}{fig:quality_tests}

\section{Commandes qualité}
\begin{lstlisting}
.\.venv\Scripts\python.exe -m pytest
.\.venv\Scripts\python.exe -m black --check src tests dags scripts
.\.venv\Scripts\python.exe -m isort --check-only src tests dags scripts
.\.venv\Scripts\python.exe -m ruff check src tests dags scripts
.\.venv\Scripts\python.exe -m mypy src
\end{lstlisting}

\section{TDD adapté au ML}
Les tests ne se limitent pas à vérifier que les fonctions retournent une valeur. Ils vérifient aussi les invariants du pipeline :
\begin{itemize}
  \item le dataset validé existe ;
  \item le modèle peut être chargé ;
  \item l'API répond ;
  \item les payloads de prédiction sont corrects ;
  \item les agrégations de l'interface SOC sont cohérentes.
\end{itemize}
